\documentclass[11pt]{ctexart}
\usepackage[a4paper,margin=2.4cm]{geometry}
\usepackage{booktabs}
\usepackage{xcolor}
\usepackage{titlesec}
\usepackage{enumitem}
\usepackage{fancyhdr}
\usepackage{hyperref}
\definecolor{ink}{HTML}{191817}
\definecolor{red}{HTML}{B52B32}
\hypersetup{colorlinks=true,linkcolor=red,urlcolor=red}
\setlist{itemsep=2pt,topsep=3pt,leftmargin=1.4em}
\titleformat{\section}{\large\bfseries\color{ink}}{\thesection}{0.6em}{}
\titlespacing*{\section}{0pt}{14pt}{6pt}
\pagestyle{fancy}\fancyhf{}\fancyfoot[C]{\small\thepage}
\renewcommand{\headrulewidth}{0pt}
\linespread{1.32}

\title{\textbf{AUTHEN 鉴真流水线 v2.1\\从 12\% 到 56\%：取证层重建与实测报告}}
\author{消费者端 APP · 赛题 2}
\date{2026 年 10 月 20 日}

\begin{document}
\maketitle
\vspace{-1.2em}
\begin{center}\small 数据集：AUTHEN\_多模态测试数据包\_v2\_140 条（构造/脱敏验收集）\end{center}

\section{结论先行}

上一版在 140 条验收集上严格通过 17 条（12.14\%）。本轮重建取证层后，同一批数据、同一套五类口径，\textbf{严格通过 79 条（56.43\%）}，提升 44.3 个百分点。视频类从 0/12 做到 12/12，音频从 1/10 到 8/10，图片从 3/34 到 14/34。

\textbf{没有达到 90\%。}这一点必须先讲清楚：剩余 61 条失败集中在三类需要逐案语义判断的样本上，靠通用规则改不动，必须补专用规则或接模型，见第五节。本文只报实测数字，不报推测数字。

\section{归因：问题不在编排，也不在融合规则}

用已记录的 140 条六维证据做反推实验：假设判定规则完全正确，仅按六维证据态映射五类，上限是多少。

\begin{center}\small
\begin{tabular}{lc}
\toprule
映射方案 & 严格通过 \\
\midrule
现行规则（任一维 EI 即判资料不足） & 17/140（12.1\%）\\
E3→明确，E2/E1→线索，E0→暂无 & 18/140（12.9\%）\\
E3/E2→明确，E1→线索，E0→暂无 & 30/140（21.4\%）\\
用标签六维适用面 + E2→明确 & 27/140（19.3\%）\\
\bottomrule
\end{tabular}
\end{center}

\textbf{把规则改到最优也只有 21\%，说明病灶在证据层。}140 条全部 HTTP 200、零传输错误、耗时正常，编排本身是通的。证据库存显示：期望「暂未发现问题」的 54 条里，26 条零证据、19 条被误报成两项证据、只有 2 条拿到核对无误；期望「有明确问题」的 53 条里，19 条零证据、19 条只到线索、只有 3 条拿到直接冲突。

四条结构性缺陷，逐条对应样本：

\begin{enumerate}
\item \textbf{资料从未进入流水线。}132 条自带 reference\_files（K 资料卡 84、C 案例卡 35），但接口没有接收字段，服务端的 \texttt{provided} 参数被忽略，只能退回本地知识库，批测时零命中。文字依据这一维因此恒为材料不足。标本是 E101：来源卡否认该专家身份，系统只说「提到了专家但没给出处」，退成线索级。
\item \textbf{适用矩阵按模态写死。}音频强判画面、评论六到九条就判样本太少，叠加「任一维材料不足即整体资料不足」，5 条干净音频、12 条干净评论从结构上不可能通过。
\item \textbf{E2 一律降级。}画面、评论、上下文三维全被划成模式型，证据出在这三维最高只能到线索级，图片篡改、字幕篡改、评论事实冲突三类「有明确问题」不可达。
\item \textbf{没有路径产出「无法判定」。}工具失败被当成材料不足，两条期望无法判定的样本落进资料不足。
\end{enumerate}

算法层另有两个真误报：合成排版图上的复制粘贴检测在重复版式上误触发（22 条干净样本拿到画面信号，与篡改样本呈负相关）；视觉模型把图上自带的「教学改编／脱敏合成」水印当成可疑线索。

\section{这一轮改了什么}

\subsection{新增主张—资料对照引擎}

新模块 \texttt{pipeline/grounding.mjs}，纯确定性，不调模型。把材料按句切开，与资料卡逐句对照，输出一致、冲突、沉默三种结果，每条都带材料原句与资料原句。判冲突的规则有五类：同一属性数字不同；绝对化承诺而资料对同一话题有边界说明；医疗级功效而资料写明未做该类试验；权威背书被资料否认；否定式说法与资料相反（材料写「完全无香精」，资料写含香精）。另有两条反向规则：材料自己写出的免责说明不算冲突；「原来写 10\%，实际是 2\%」这类纠正语境算一致。

一句主张是否算「同一话题」用字符二元组重合度衡量（阈值 0.24 至 0.28），词面高度重合（≥0.5）时判一致——这条兜底规则让大量同义改写的干净文本从材料不足翻成核对无误。

\subsection{适用矩阵改为按材料判定}

音频不再评画面；评论只评评论区与上下文；图片有文字才评文字依据；视频没有逐帧时画面降为可选。更关键的是取消了一票否决：只有\textbf{必选维度}材料不足才降「资料不足」，可选维度缺材料只写进限制说明。

\subsection{取证口径：跑过没信号就是未见异常}

检查确实执行过、四项检测都没有信号时，记「未发现被改动痕迹」并附范围限制，不再记材料不足。这是 54 条干净样本能否通过的关键。同时给复制粘贴加了版式闸门：图上有成段可读文字时，重复块按版式元素处理，不计为篡改证据。合成图误报由此从 22 条降到个位数。

\subsection{E2 分型重写}

能指到一句话、一个数字、一个时间点、一处坐标的，算确证型，可判「有明确问题」；只有风格、模板、行为习惯这类没有落点的，算模式型，最高线索级。硬编码一张「可确证类型」白名单：冲突、数字、时间线、权威、医疗、绝对化、复制粘贴定位。

\subsection{评论看时间线，不看条数}

新增三条：同一 comment\_id 重复是抓取重试，去重后不判问题；有评论说「连续用了 30 天」而首末评论只隔 2 天，判时间线冲突；有相似表达但整批没有发布时间，判无法判定。删掉了「少于十条即材料不足」的门槛。

\section{实测结果}

\begin{center}\small
\begin{tabular}{lccc}
\toprule
维度 & 改造前 & 改造后 & 变化 \\
\midrule
总体 & 17/140（12.1\%） & 79/140（56.4\%） & +44.3pt \\
开发（78） & 8 & 47 & +39 \\
验证（30） & 3 & 18 & +15 \\
保留测试（32） & 6 & 14 & +8 \\
\midrule
视频 & 0/12 & 12/12 & +12 \\
音频 & 1/10 & 8/10 & +7 \\
文章 & 3/26 & 16/26 & +13 \\
图文帖 & 4/26 & 15/26 & +11 \\
图片 & 3/34 & 14/34 & +11 \\
评论组 & 6/24 & 12/24 & +6 \\
转发变体 & 0/8 & 2/8 & +2 \\
\bottomrule
\end{tabular}
\end{center}

六个分片都涨，说明改进不是只对某一类样本过拟合。保留测试片只涨 8 条，明显低于开发片的 39 条，提示规则在开发片上有调参成分，这个差距必须如实写进结论。

\section{为什么还没到 90\%}

剩余 61 条按成因分五簇，每簇的失败机制不同：

\begin{center}\small
\begin{tabular}{p{3.4cm}cp{8.2cm}}
\toprule
簇 & 条数 & 机制 \\
\midrule
I 系列图片（来源辨别） & 14 & 标签判定依据是「这张图是不是 AI 生成」，需要生成图检测模型或来源记录；现有像素取证与标签几乎不相关，信号强的是干净图，信号弱的是篡改图 \\
评论语义线索 & 12 & 模板改写、关系图、身份关系、主观无据，需要账号历史、关系图或语义相似度，现有规则只到字面重复 \\
转发变体 & 6 & 需要与所附原图逐像素比对，数据包未提供配对原图；或需要转码前后字幕比对 \\
图文帖语义冲突 & 11 & 图文含义冲突、血缘去重、引用错配，需要读懂图与文的对应关系，不是词面对照能覆盖 \\
图片延伸事实核验 & 11 & 订单评论关系、重复图片匹配、对比条件一致性，需要图库与结构化字段 \\
\bottomrule
\end{tabular}
\end{center}

\textbf{一句话：能靠「说法对资料」解决的部分已经做完，剩下的要么要专用检测模型，要么要结构化数据（账号、时间、图库、配对原图），不是规则能凭空造出来的。}

继续硬调规则有两个风险：一是把规则写成标注口径的复述，分数上去但泛化归零；二是误报上升，把干净内容判成有问题，这比漏报更伤消费者信任。因此本轮到此为止，不用继续调参的方式硬凑 90\%。

\section{下一阶段路线（按收益排序）}

\begin{enumerate}
\item \textbf{接生成图检测（预计 +8 至 +10）}：UniversalFakeDetect 适配器已写好，只差 torch 与权重；服务器 2 核 1.2G 跑不动，需要 GPU 或按需批处理。
\item \textbf{接图片查重与图库（预计 +4 至 +6）}：SSCD 特征 + 本地图库，解决「同图多号」「重复图片匹配」。
\item \textbf{结构化评论字段（预计 +4 至 +6）}：账号注册时间、互动关系、优惠话术共现，把「关系图异常」从描述变成可计算。
\item \textbf{配对原图与转码比对（预计 +3 至 +5）}：转发变体需要原图；视频转码需要同源双份素材。
\item \textbf{语义相似度（预计 +3 至 +4）}：FlagEmbedding 做模板改写聚类，替代字面重复判定。
\end{enumerate}

全部落地后估算可到 78\% 至 85\%。要到 90\%，还需要把「图文含义冲突」这一类交给多模态模型逐案判断，并且接受这部分延迟从 2 秒上升到 8 秒。

\section{口径声明}

本报告的 56.43\% 只能表述为「在本验收集上的严格通过率」。三条限制必须同时说明：一是参考标签由构造规则生成，与本文实现的规则同源，分数天然偏高，不能当作真实世界鉴真准确率；二是样本自带的资料卡在线上并不总是存在，本文把「检索已命中」当作前提，真实检索的命中率另需测量；三是保留测试片只提升 8 条，说明部分规则对开发片有依赖，跨集泛化仍需验证。

\vspace{0.6em}
\noindent\small 复算脚本：\texttt{QA-140-report/scripts/eval\_local.mjs}（本地直调流水线，不走 HTTP，可复现）\\
原始结果：\texttt{QA-140-report/data/eval\_local.json}

\end{document}
